\begin{center}
\begin{tikzpicture}[x=.88cm,y=.88cm,font=\small]
  \fill[paper] (0,0) rectangle (12.3,6.65);

  \draw[source!70,line width=.8pt] (.52,2.0) .. controls (3.4,1.65) and
    (7.05,2.54) .. (11.85,1.78);
  \node[text=source,below,font=\scriptsize] at (8.28,1.84) {长江方向（示意）};
  \draw[source!55,dashed] (.78,4.54) .. controls (3.7,4.92) and
    (7.45,4.04) .. (11.7,4.48);
  \node[text=source,above,font=\scriptsize] at (8.82,4.3) {淮河方向（示意）};

  \fill[dynasty!44] (4.0,4.95) ellipse (1.5 and .92);
  \node[align=center,text=white] at (4.0,4.95) {隋\\大兴、关中中枢};

  \fill[timeline!25] (1.35,2.7) ellipse (1.08 and .75);
  \node[align=center,font=\scriptsize] at (1.35,2.7) {江陵\\后梁旧地};

  \fill[source!25] (6.88,3.55) ellipse (1.38 and .86);
  \node[align=center] at (6.88,3.55) {荆襄、淮南\\渡江前进节点};

  \fill[dynasty!34] (10.18,1.0) ellipse (1.48 and .76);
  \node[align=center,text=white] at (10.18,1.0) {陈\\建康与江东};

  \fill[timeline!21] (9.98,4.92) ellipse (1.13 and .75);
  \node[align=center,font=\scriptsize] at (9.98,4.92) {淮北、江北\\隋军诸路};

  \draw[->,very thick,color=dynasty]
    (4.94,4.55) .. controls (5.45,4.16) and (5.84,3.88) .. (6.03,3.77);
  \node[text=dynasty,above,font=\scriptsize,align=center] at (5.55,4.56)
    {户籍、粮道与\\分路动员};

  \draw[->,thick,dashed,color=timeline]
    (2.26,3.1) .. controls (3.42,3.54) and (4.67,3.66) .. (5.43,3.63);
  \node[text=timeline,below,font=\scriptsize,align=center] at (3.72,2.95)
    {587年废后梁后\\取得荆襄道路};

  \draw[->,very thick,color=dynasty]
    (7.9,3.23) .. controls (8.64,2.62) and (9.2,1.92) .. (9.48,1.54);
  \node[text=dynasty,left,font=\scriptsize,align=right] at (8.68,2.4)
    {588--589年\\向建康推进};

  \draw[->,thick,dashed,color=source]
    (9.68,4.22) .. controls (9.35,3.53) and (9.14,3.12) .. (9.02,2.85);
  \node[text=source,right,font=\scriptsize,align=left] at (9.82,3.42)
    {江北诸军\\并进与牵制};

  \node[draw=source!75,fill=paper,rounded corners=2pt,align=center,
    inner sep=3pt,font=\scriptsize] at (3.08,.72)
    {589年陈亡解决最高政权归属；州郡、户籍、仓储、官员和江南社会\\
    的接收仍是渡江之后的工作};
\end{tikzpicture}

\smallskip
\small\textit{图\quad 隋陈再统一前后的关中、荆襄、江淮与建康，约581--589年：据《隋书·高祖纪下》《陈书·后主纪》、〈杨素传〉〈韩擒虎传〉及《资治通鉴》卷一百七十五至一百七十七，概括隋的北方中枢、废后梁后取得的江陵道路与伐陈诸路。节点位置和河流为约略示意；箭线不等于军队的实际渡江、行军或战斗路线，色块也不表示可精确测量的控制区。}
\end{center}

\phantomsection\label{fig:sui-chen-reunification}
